% Level 2, Week 5 -- Dictionaries and Secret Codes
\lesson{5}{Dictionaries and Secret Codes}{A dictionary looks things up by name, not by number, and can even hold a secret code.}
\lessoninfo{Concepts: dict, keys/values, .items() (I), Caesar cipher (I), strings (R)}{Tools: Thonny, paper cipher wheel, scissors, split pin or pencil}{CSTA: 2-AP-11, 3A-AP-14, 2-NI-05, 2-NI-06}

\begin{objectives}
  \item The student can make a dictionary, read, change and add keys, and loop over it with \texttt{.items()}. (L2.G4; L2.W05.1)
  \item The student can encode and decode a message with a Caesar shift, on paper and in Python. (L2.G4; L2.W05.2)
  \item The student can explain what a \texttt{KeyError} means and avoid it with \texttt{if key in dictionary}. (L2.G6; L2.W05.3)
  \item The student can explain why passwords and other secrets need protecting, and never type a real password into our programs. (L2.G9; L2.W05.4)
\end{objectives}

\begin{materials}
  \item Print (one per student): workbook chapter 5 (it contains the cipher wheel page), homework 5.
  \item Scissors, and a split pin (paper fastener) or a pencil point to hold the two wheels together.
  \item Before the session (5 minutes): write a secret message on a sticky note with shift 3: \texttt{brx duh d frgh euhdnhu}
        (it says \enquote{you are a code breaker}).
  \item On each Mac: in Thonny press \keys{Cmd}+\keys{O} and open \texttt{starter.py} in the folder \texttt{2-intermediate/} \texttt{excercises/week05/ex-monster-stats}.
  \item A timer.
\end{materials}

\needspace{16\baselineskip}
\section{At a glance}
Times are a guide: set a timer for each phase and move on when it rings.
\begin{timeplan}
  \timeslot{0--8}{Warm-up}{Unplugged: build the paper Caesar wheel and crack your secret sticky note.}
  \timeslot{8--18}{Concept}{Lists find things by number; dictionaries find things by name. Keys and values. KeyError.}
  \timeslot{18--28}{Guided demo}{Build a monster dictionary in the Shell; loop with \texttt{.items()}; trigger and fix a \texttt{KeyError}.}
  \timeslot{28--52}{Hands-on build}{Pairs: Exercise 5.1 Monster Stats, then 5.2 Caesar Cipher. Swap every 8 minutes. Last 5 minutes: secret-message swap.}
  \timeslot{52--60}{Wrap-up}{Encryption and passwords talk, exit ticket, reflection, hand out homework 5.}
\end{timeplan}

\section{Warm-up (0--8 min)}
\begin{unplugged}[Paper Caesar wheel]
Students cut out the two wheels from the workbook page, put the small one on top of the big one and pin the centres
(split pin, or a pencil point). Turn the inner wheel so that inner \textbf{a} sits under outer \textbf{d}: that is a \textbf{shift of 3}.
\par To \textbf{encode}, find the plain letter on the inner wheel and write the outer letter. To \textbf{decode}, go the other way: outer to inner.
\par Show your sticky note \texttt{brx duh d frgh euhdnhu} and say: \enquote{This was made with shift 3. Who can crack it?}
(Answer: \enquote{you are a code breaker}.) Spaces stay spaces.
\par\textbf{Ask:} \enquote{What comes after z?} (The wheel goes round, so z becomes c with shift 3.) Keep the wheels for homework and for B1.
\end{unplugged}

\section{Concept (8--18 min)}
\begin{teacherscript}[Say]
\enquote{Your cipher wheel is a lookup: give it a letter, it gives you the secret letter. Python has a tool made for lookups:
the \textbf{dictionary}. A list finds things by \emph{number}: \texttt{backpack[0]}. A dictionary finds things by \emph{name}.
Each name is called a \textbf{key}, and what it gives back is the \textbf{value}. We write a dictionary with curly brackets,
and a colon between each key and its value.}
\end{teacherscript}
Write on paper, with arrows from each key to its value:
\begin{pycode}
monster = {"name": "Blobby", "health": 30, "attack": 4}
monster["health"]              # 30   look up by key
monster["health"] = 40         # change a value
monster["colour"] = "green"    # add a new key
monster["magic"]               # KeyError: 'magic'  (no such key!)
\end{pycode}

\begin{analogy}[The real dictionary]
In a paper dictionary you do not read page 1, page 2, page 3\ldots{} to find \enquote{dragon}: you jump straight to the word and read its meaning.
The word is the \textbf{key}; the meaning is the \textbf{value}. If you look up a word that is not in the book, there is nothing to find:
that is Python's \texttt{KeyError}. And like a real dictionary, the spelling must be exact: \texttt{"Health"} and \texttt{"health"} are different words.
\end{analogy}

\section{Guided demo (18--28 min)}
Type each line below into the Thonny \textbf{Shell}, pressing \keys{Return} after each. Ask for a prediction first.
\begin{pycode}[Monster stats in the Shell]
blobby = {"name": "Blobby", "health": 30, "attack": 4}
blobby["attack"]                          # 4
blobby["attack"] = blobby["attack"] + 2   # level up the attack
blobby                                    # see the whole dictionary
for key, value in blobby.items():
    print(f"{key}: {value}")              # press Return twice to run the loop
blobby["Attack"]                          # KeyError (capital A!)
"magic" in blobby                         # False: check first
\end{pycode}
\begin{workedexample}[Talking through the Shell demo]
\begin{enumerate}[leftmargin=1.4em]
  \item After the \texttt{for} line, say: \enquote{\texttt{.items()} hands us each key \emph{and} its value, one pair at a time.}
  \item After the \texttt{KeyError}, read it bottom-up: \enquote{KeyError, and then the key it could not find, in quotes. Look at the spelling.}
  \item Then open \texttt{ex-monster-stats/starter.py}, press \menu{Run} (\keys{F5}) and type \texttt{magic} when it asks for a stat. Crash!
        That fix (TODO 3) is for the students.
\end{enumerate}
\end{workedexample}
The finished Monster Stats, for you:
\pyfile{../excercises/week05/ex-monster-stats/solution.py}

\section{Hands-on build (28--52 min)}
\begin{pairroles}[Codemakers]
Swap driver and navigator when the 8-minute timer rings (about minutes 36 and 44).
\textbf{1 student:} you are the navigator; in the last 5 minutes you send them a secret message to decode with their program.
\textbf{2 students:} one pair. \textbf{3 students:} a rotating trio with a \textbf{codebreaker} who checks every encoded message with the paper wheel.
\textbf{4 students:} two pairs. In the free-build, each pair encodes a message for the other pair.
\end{pairroles}

\begin{enumerate}[leftmargin=1.4em]
  \item \textbf{Exercise 5.1 Monster Stats} (folder \texttt{ex-monster-stats}, about 10 minutes): TODO 1 (\texttt{.items()}), TODO 2 (level up), TODO 3 (KeyError).
  \item \textbf{Exercise 5.2 Caesar Cipher} (folder \texttt{ex-caesar-cipher}, about 12 minutes). Run it first and look at the printed dictionary
        \texttt{\{'a': 'd', 'b': 'e', ...\}}: \enquote{That is your paper wheel, as a dictionary!} Then finish \texttt{encode()} and \texttt{decode()}.
  \item After each exercise: open \texttt{check\_my\_work.py} and press \menu{Run}.
\end{enumerate}

\begin{teachernote}[The \% sign in make\_key()]
\texttt{(i + shift) \% 26} is the remainder after dividing by 26: it wraps the alphabet round like a clock, so after \textbf{z} comes \textbf{a}.
Students only use \texttt{make\_key()}; if asked, say \enquote{it is your paper wheel going round, in code}.
\end{teachernote}

\begin{bughunt}[KeyError and wrong spelling]
Two kinds of bug today. (1) The loud one: \texttt{get\_stat(blobby, "magic")} crashes with \texttt{KeyError: 'magic'}. Fix: \texttt{if stat in monster:}
before looking it up, and \texttt{return 0} otherwise. (2) The quiet one: \texttt{monster["Health"] = 50} does \emph{not} crash. It adds a new
key with a capital H, and the real \texttt{"health"} does not change. Show it in the Shell and ask: \enquote{Why are there two healths now?}
Printing the whole dictionary is the fastest way to catch it.
\end{bughunt}

\begin{freebuild}[Secret-message swap]
Last 5 minutes: each student (or pair) uses their Caesar program to encode a short, kind message with a secret shift,
writes the secret version on a sticky note, and passes it on. The receiver decodes it, with the program or the paper wheel.
Extension: decode one without being told the shift.
\end{freebuild}

\section{Wrap-up (52--60 min)}
\begin{safety}[Encryption and passwords]
Say: \enquote{You just cracked a Caesar cipher in a minute! Real apps use much stronger \textbf{encryption}, maths so hard
that even fast computers cannot undo it without the key. That keeps messages and bank details safe.} Agree these rules together:
(1) a real password is only for you and your parent or carer: not for friends, games or chat;
(2) never type a real password into a program you or a friend wrote: we only use made-up game passwords;
(3) long is strong: three random words joined together beat a short word with a number;
(4) if a website or message asks for your password, ask a grown-up first: real companies never ask for it in a message.
\end{safety}

\begin{formativecheck}[Exit ticket]
\begin{pycode}
pet = {"name": "Pixel", "hunger": 5}
pet["hunger"] = pet["hunger"] - 2
print(pet["hunger"])
print(pet["happy"])
\end{pycode}
What happens on lines 3 and 4? \textbf{Answer:} line 3 prints \texttt{3}; line 4 crashes with \texttt{KeyError: 'happy'} (no such key).
Bonus: encode \texttt{cat} with shift 1 (\texttt{dbu}). \textbf{Look for:} mixing up keys and values (\enquote{it prints Pixel}): draw the arrows again.
\end{formativecheck}

\begin{reflection}[Say it or write it]
\begin{itemize}[leftmargin=1.2em]
  \item When is a dictionary better than a list? Give a game example.
  \item What does a \texttt{KeyError} tell you, and how do you fix it?
  \item Why is a Caesar cipher not good enough for real secrets?
\end{itemize}
\end{reflection}
Hand out homework 5 (Word Counter). It takes 15--20 minutes.

\section{Notes for the teacher}
\begin{commonmistakes}
  \item \textbf{Square or round brackets instead of curly ones} when making a dictionary. Making: \texttt{\{ \}}. Looking up: \texttt{monster["health"]}.
  \item \textbf{Key without quotes} (\texttt{monster[health]} gives \texttt{NameError}), or \textbf{wrong spelling/capitals} (a quiet bug that makes a new key: print the dictionary to spot it).
  \item \textbf{\texttt{for key in monster:}} gives only the keys. To get both, use \texttt{for key, value in monster.items():}.
  \item \textbf{Caesar: losing the spaces}. Letters not in the key (spaces, \texttt{!}) must be added unchanged in an \texttt{else}.
  \item \textbf{Decode shifts the wrong way}: decoding with shift 3 is encoding with shift \texttt{-3} (or 23).
\end{commonmistakes}

\begin{differentiation}
\supportpath Draw the monster dictionary as a key/value table on paper before coding. In 5.2, do \texttt{encode()} together on paper
for \texttt{hi}: look up h, add the result, then i. Hint for \texttt{decode()}: \enquote{it is one line}.
\extensionpath Write \texttt{crack(secret)} that prints all 25 shifts so you can spot the real message (try it on \texttt{brx fudfnhg lw}). In 5.1, write
\texttt{battle(a, b)} that returns the name of the monster with the higher attack, and store several monsters in a list of dictionaries.
\end{differentiation}

\begin{stretchchallenge}[Keep the capitals]
Our cipher turns every message into small letters. Make \texttt{encode()} keep capital letters: if \texttt{letter.isupper()}, look up
\texttt{letter.lower()} and add the result with \texttt{.upper()}. Check that \texttt{Hi Zoe!} with shift 3 becomes \texttt{Kl Crh!}.
\end{stretchchallenge}

\begin{teachernote}[If you are short on time]
Skip building the wheels in class: decode the sticky note together with the alphabet strip in the workbook instead, and send the wheel home.
Do 5.2 Caesar Cipher in class and set 5.1 Monster Stats as an extra. Keep the passwords talk: it is short and important.
\end{teachernote}
